\chapter{Nền tảng công nghệ}
\textit{Chương này phân tích lý do lựa chọn các công nghệ cốt lõi cho nền tảng ôn luyện tiếng Anh tích hợp AI. Mỗi quyết định được đánh giá dựa trên các yêu cầu đặc thù của dự án: xử lý nội dung giáo dục đa phương tiện, tích hợp AI inference, khả năng mở rộng cho lượng người dùng lớn, và tốc độ phát triển phù hợp với tiến độ đồ án.}

\section{Công nghệ Frontend}

\subsection{Phân tích và lựa chọn frontend framework}

Nền tảng ôn luyện tiếng Anh có ba yêu cầu frontend đặc thù: Thứ nhất là SEO tốt để người dùng tìm thấy nội dung luyện thi qua search engine, hai là tối ưu hiệu năng tải trang cho nội dung giàu hình ảnh và audio, và cuối cùng là hỗ trợ việc truyền tải dữ liệu, trang tĩnh (giới thiệu, blog) lẫn trang động (dashboard, làm bài thi). Ba framework được đánh giá dựa trên các tiêu chí này:

\textbf{AngularJS} sử dụng kiến trúc MVC với two-way data binding, nhưng không hỗ trợ server-side rendering (SSR) mặc định, gây hạn chế nghiêm trọng đối với yêu cầu SEO của nền tảng giáo dục. Bundle size lớn do tính chất monolithic cũng ảnh hưởng đến thời gian tải ban đầu. \cite{angularjs}

\textbf{ReactJS} Framework mạnh về xây dựng UI tương tác nhờ Virtual DOM và kiến trúc component-based, nhưng client-side rendering (CSR) mặc định gây bất lợi cho SEO. Cần bổ sung nhiều thư viện riêng cho routing, quản lý state và tối ưu việc truyền tải hình ảnh. \cite{angular_react}

\textbf{Next.js} Được xây dựng trên React, bổ sung hybrid rendering (SSR + SSG), code splitting tự động và có image optimization giải quyết trực tiếp cả ba yêu cầu đặc thù của dự án mà không cần cấu hình thêm. \cite{next_react_angular}

\begin{table}[H]
\centering
\resizebox{\textwidth}{!}{%
\begin{tabular}{|l|p{4cm}|p{4cm}|p{4cm}|}
\hline
\textbf{Tiêu chí dự án} & \textbf{AngularJS} & \textbf{ReactJS} & \textbf{Next.js} \\ \hline
Hỗ trợ SEO cho nội dung giáo dục & Kém, cần Angular Universal riêng & Trung bình, render phía client, cần cấu hình SSR thủ công & Tích hợp sẵn SSR/SSG, index tốt bởi search engine \\ \hline
Tối ưu nội dung đa phương tiện & Cần thư viện bên thứ ba & Cần thư viện bên thứ ba & Tích hợp sẵn image optimization, tự động chuyển WebP \\ \hline
Hỗ trợ hybrid (tĩnh + động) & Không hỗ trợ native & Cần thiết lập thủ công & SSG cho trang tĩnh, SSR cho dashboard cùng một app \\ \hline
Tốc độ phát triển & Chậm,learning curve dốc, boilerplate nhiều & Nhanh cho UI, nhưng cần tích hợp nhiều library & Nhanh, routing, API routes, TypeScript đều tích hợp sẵn \\ \hline
\end{tabular}%
}
\caption{So sánh frontend framework theo yêu cầu dự án}
\end{table}

Nhóm chọn \textbf{Next.js} vì framework này đáp ứng trực tiếp các yêu cầu đặc thù của nền tảng giáo dục:

\begin{enumerate}
\item \textbf{SEO cho nội dung luyện thi}: Nền tảng phụ thuộc vào nguồn tìm kiếm tự nhiên để tiếp cận học viên. SSR tích hợp sẵn bảo đảm nội dung bài thi, bài viết hướng dẫn được search engine index đầy đủ, là yếu tố quyết định khả năng tăng trưởng người dùng mà không cần ngân sách quảng cáo lớn.

\item \textbf{Tối ưu tài liệu giàu hình ảnh}: Tài liệu luyện thi chứa nhiều hình ảnh minh họa, sơ đồ và biểu đồ. Next.js Image tự động phân phối hình ảnh đúng kích thước ở định dạng WebP, giảm đáng kể thời gian tải, đặc biệt quan trọng cho người dùng ở vùng có kết nối Internet chậm.

\item \textbf{Hybrid rendering phù hợp kiến trúc}: Trang giới thiệu và blog dùng SSG (tải nhanh, cache tốt), trong khi dashboard làm bài và kết quả AI dùng SSR (dữ liệu luôn cập nhật). Next.js cho phép kết hợp cả hai trong cùng một ứng dụng mà không cần tách dự án.

\item \textbf{TypeScript xuyên suốt}: Sử dụng chung TypeScript với backend NestJS, giúp nhóm tái sử dụng type definition (DTO, API response) giữa frontend và backend, giảm lỗi runtime và tăng tốc phát triển.
\end{enumerate}

\section{Công nghệ Backend}
\subsection{Phân tích và lựa chọn backend framework}

Hệ thống cần một backend framework đáp ứng các yêu cầu: application logic hướng người dùng (authentication, CRUD, real-time notification), kiến trúc microservices, và khả năng tích hợp với các dịch vụ AI thông qua hàng đợi bất đồng bộ.

\textbf{NestJS} (TypeScript/Node.js): framework modular với dependency injection, lấy cảm hứng từ Angular. Mô hình event-driven, non-blocking I/O của Node.js phù hợp cho workload concurrency cao (WebSocket, real-time dashboard). Chia sẻ TypeScript với frontend Next.js giúp tái sử dụng type definition và giảm context switching cho nhóm phát triển. \cite{nestjs}

\textbf{ASP.NET Core} (C\#) và \textbf{Spring Boot} (Java): đều là framework enterprise-grade với performance và reliability cao. Tuy nhiên, cả hai đều có chu kỳ phát triển dài hơn, không phù hợp với tiến độ ngắn của đồ án. \cite{aspnet, springboot}

\textbf{FastAPI} (Python): framework API hiệu năng cao với async native. Phù hợp cho AI/ML workload nhưng không phải thế mạnh cho application logic phức tạp (auth, CRUD, real-time). \cite{fastapi}

\begin{table}[H]
    \centering
    \caption{So sánh backend framework theo yêu cầu dự án}
    \label{tab:framework_comparison}
    \scriptsize
    \resizebox{\textwidth}{!}{%
    \begin{tabular}{|p{2.8cm}|p{3.5cm}|p{3.5cm}|p{3.5cm}|}
        \hline
        \textbf{Tiêu chí dự án} & \textbf{NestJS (TypeScript)} & \textbf{ASP.NET / Spring Boot} & \textbf{FastAPI (Python)} \\
        \hline
        Application logic (auth, CRUD, real-time) & Rất mạnh, event-driven, modular, WebSocket native & Mạnh nhưng chu kỳ phát triển dài hơn & Khả dụng nhưng không phải thế mạnh \\
        \hline
        Kiến trúc Microservices & Hỗ trợ tốt qua module system, CQRS, message queue & Hỗ trợ tốt nhưng cấu hình phức tạp & Hỗ trợ cơ bản \\
        \hline
        Tốc độ phát triển (tiến độ đồ án) & Rất nhanh, chia sẻ TypeScript với frontend & Chậm hơn, boilerplate nhiều & Nhanh cho API đơn giản \\
        \hline
        Tái sử dụng code với frontend & Có, chung TypeScript type definitions & Không & Không (Python vs TypeScript) \\
        \hline
    \end{tabular}
    }
\end{table}

\textbf{Quyết định}: Sử dụng \textbf{NestJS} làm backend framework duy nhất cho toàn bộ hệ thống, với các lý do:

\begin{itemize}
\item \textbf{Kiến trúc modular:} NestJS cho phép tổ chức hệ thống thành các module độc lập (auth, exam, achievement, file-service), dễ phân chia công việc trong nhóm 4 người và phù hợp với kiến trúc microservices.
\item \textbf{Chia sẻ TypeScript với frontend:} Cùng ngôn ngữ TypeScript với Next.js frontend, cho phép tái sử dụng DTO và API response type, giảm lỗi tích hợp và context switching.
\item \textbf{Hỗ trợ real-time và message queue:} Event-driven architecture, WebSocket native và tích hợp RabbitMQ giúp xử lý thông báo thời gian thực và giao tiếp bất đồng bộ với AI server.
\item \textbf{Tích hợp AI qua queue:} Các tác vụ AI inference (chấm bài Writing, phân tích Speaking) được giao tiếp qua RabbitMQ với AI server riêng biệt, không cần tích hợp AI library trực tiếp vào backend.
\end{itemize}

\subsection{Phân tích và lựa chọn database}

Hệ thống cần một hệ quản trị cơ sở dữ liệu đáp ứng các yêu cầu: ACID compliance cho dữ liệu transactional (user, enrollment, permission), hỗ trợ quan hệ phức tạp giữa các thực thể, và khả năng lưu trữ dữ liệu bán cấu trúc khi cần.

Bốn hệ quản trị database được đánh giá:

\begin{table}[H]
    \centering
    \caption{So sánh database theo yêu cầu dự án}
    \label{tab:db_comparison}
    \scriptsize
    \resizebox{\textwidth}{!}{%
    \begin{tabular}{|p{2.2cm}|p{2cm}|p{4.5cm}|p{4.5cm}|}
        \hline
        \textbf{Database} & \textbf{Loại} & \textbf{Phù hợp cho dự án} & \textbf{Hạn chế với dự án} \\
        \hline
        \textbf{PostgreSQL} & Relational & ACID mạnh cho auth/enrollment; complex query cho báo cáo; JSONB cho dữ liệu bán cấu trúc & Schema migration cần quản lý cẩn thận khi format bài thi thay đổi \\
        \hline
        \textbf{MongoDB} & Document NoSQL & Schema linh hoạt cho nội dung giáo dục; horizontal scaling & Không bảo đảm ACID đầy đủ cho transaction phức tạp; thêm công nghệ vào stack \\
        \hline
        \textbf{MySQL} & Relational & Đơn giản, cộng đồng lớn & Thiếu tính năng nâng cao (JSONB, full-text search) so với PostgreSQL \\
        \hline
        \textbf{SQL Server} & Relational & Enterprise-grade, BI tốt & Chi phí license cao, thiên về Windows ecosystem \cite{sqlserver} \\
        \hline
    \end{tabular}
    }
\end{table}

\textbf{Quyết định}: Sử dụng \textbf{PostgreSQL} làm hệ quản trị cơ sở dữ liệu duy nhất cho toàn bộ hệ thống, với các lý do:

\begin{itemize}
\item Các workflow transactional quan trọng (authentication, ghi danh, quản trị, chấm điểm) đòi hỏi ACID compliance nghiêm ngặt mà PostgreSQL bảo đảm đáng tin cậy.
\item Relational model phù hợp cho dữ liệu có quan hệ rõ ràng: user $\rightarrow$ enrollment $\rightarrow$ exam result $\rightarrow$ achievement.
\item Hỗ trợ JSONB cho phép lưu metadata linh hoạt (cấu trúc câu hỏi thi, kết quả phân tích AI) mà không cần thêm một hệ quản trị NoSQL riêng.
\item Giảm độ phức tạp vận hành: chỉ cần quản lý một loại database, đơn giản hóa backup, monitoring và deployment. \cite{postgresql}
\end{itemize}

\subsection{Phân tích và lựa chọn giải pháp lưu trữ đối tượng}

Hệ thống cần một giải pháp lưu trữ đối tượng (object storage) để quản lý các file đa phương tiện như audio ghi âm Speaking, hình ảnh bài thi, tài liệu đính kèm và các file kết quả phân tích AI. Ba giải pháp được đánh giá dựa trên yêu cầu dự án:

\begin{table}[H]
    \centering
    \caption{So sánh giải pháp lưu trữ đối tượng theo yêu cầu dự án}
    \label{tab:storage_comp}
    \scriptsize
    \resizebox{\textwidth}{!}{%
    \begin{tabular}{|p{3cm}|p{4cm}|p{4cm}|p{3.5cm}|}
        \hline
        \textbf{Tiêu chí dự án} & \textbf{MinIO} & \textbf{Amazon S3} & \textbf{Google Cloud Storage} \\
        \hline
        Chi phí & Mã nguồn mở, miễn phí khi tự triển khai & Trả phí theo dung lượng và lượt truy cập & Trả phí theo dung lượng và lượt truy cập \\
        \hline
        Tương thích S3 API & Hoàn toàn tương thích, có thể chuyển sang S3 khi cần & Native & Cần adapter riêng \\
        \hline
        Triển khai trên VPS & Dễ dàng qua Docker Compose & Không hỗ trợ, chỉ dùng trên AWS & Không hỗ trợ, chỉ dùng trên GCP \\
        \hline
        Kiểm soát dữ liệu & Toàn quyền, dữ liệu lưu trên máy chủ của nhóm & Dữ liệu trên hạ tầng AWS & Dữ liệu trên hạ tầng Google \\
        \hline
    \end{tabular}
    }
\end{table}

MinIO \cite{minio} được chọn vì ba lý do phù hợp với bối cảnh dự án. Thứ nhất, MinIO là giải pháp mã nguồn mở có thể triển khai trực tiếp trên VPS thông qua Docker Compose, không phát sinh chi phí lưu trữ đám mây. Thứ hai, MinIO tương thích hoàn toàn với S3 API, cho phép nhóm sử dụng các thư viện S3 SDK sẵn có trong cả NestJS và FastAPI, đồng thời dễ dàng chuyển sang Amazon S3 khi cần mở rộng mà không phải thay đổi mã nguồn. Thứ ba, MinIO Console cung cấp giao diện quản trị trực quan để giám sát dung lượng, quản lý bucket và theo dõi hoạt động tải file, phù hợp cho giai đoạn phát triển và vận hành thử nghiệm.

